%% a file containing just the text macros from Brockie's typesetting file


%% colours

\newcommand{\red}[1]{\textcolor{red}{#1}}
\newcommand{\blue}[1]{\textcolor{blue}{#1}}
\newcommand{\green}[1]{\textcolor{green}{#1}}

%% Bold greek quantities
\newcommand{\bxi}{\mbox{\boldmath$\xi$}}

%% Defining aliases...

\newcommand{\ie}{{\em i.e. }}
\newcommand{\eg}{{\em e.g. }}
\newcommand{\etal}{{\em et al. }}

\newcommand{\be}{\begin{equation}}
\newcommand{\ee}{\end{equation}}
\newcommand{\ba}{\begin{eqnarray}}
\newcommand{\ea}{\end{eqnarray}}

\newcommand{\glos}[2]{
\begin{tabular}{ll}
	\parbox[t]{0.16 \textwidth}{#1} &
	\parbox[t]{0.76 \textwidth}{#2} \\
\end{tabular}
}

\newcommand{\apriori}{{\it a priori} }

%%ADS abstract ervice Bibtex entries

\newcommand{\aj}{Astronomical Journal}
\newcommand{\araa}{Annual Review of Astronomy and Astrophysics}
\newcommand{\apj}{Astrophysical Journal}
\newcommand{\apjl}{Astrophysical Journal, Letters}
\newcommand{\apjs}{Astrophysical Journal, Supplement}
\newcommand{\ao}{Applied Optics}
\newcommand{\apss}{Astrophysics and Space Science}
\newcommand{\aap}{Astronomy and Astrophysics}
\newcommand{\aapr}{Astronomy and Astrophysics Reviews}
\newcommand{\aaps}{Astronomy and Astrophysics, Supplement}
\newcommand{\azh}{Astronomicheskii Zhurnal}
\newcommand{\baas}{Bulletin of the AAS}
\newcommand{\jrasc}{Journal of the Royal Astronomical Society of Canada}
\newcommand{\memras}{Memoirs of the Royal Astronomical Society}
\newcommand{\mnras}{Monthly Notices of the Royal Astronomical Society}
\newcommand{\pra}{Physical Review A: General Physics}
\newcommand{\prb}{Physical Review B: Solid State}
\newcommand{\prc}{Physical Review C}
\newcommand{\prd}{Physical Review D}
\newcommand{\pre}{Physical Review E}
\newcommand{\prl}{Physical Review Letters}
\newcommand{\pasp}{Publications of the ASP}
\newcommand{\pasj}{Publications of the ASJ}
\newcommand{\qjras}{Quarterly Journal of the Royal Astronomical Society}
\newcommand{\skytel}{Sky and Telescope}
\newcommand{\solphys}{Solar Physics}
\newcommand{\sovast}{Soviet Astronomy}
\newcommand{\ssr}{Space Science Reviews}
\newcommand{\zap}{Zeitschrift fuer Astrophysik}
\newcommand{\nat}{Nature}
\newcommand{\iaucirc}{IAU Cirulars}
\newcommand{\aplett}{Astrophysics Letters}
\newcommand{\apspr}{Astrophysics Space Physics Research}
\newcommand{\bain}{Bulletin Astronomical Institute of the Netherlands}
\newcommand{\fcp}{Fundamental Cosmic Physics}
\newcommand{\gca}{Geochimica Cosmochimica Acta}
\newcommand{\grl}{Geophysics Research Letters}
\newcommand{\jcp}{Journal of Chemical Physics}
\newcommand{\jgr}{Journal of Geophysics Research}
\newcommand{\jqsrt}{Journal of Quantitiative Spectroscopy and Radiative Transfer}
\newcommand{\memsai}{Mem. Societa Astronomica Italiana}
\newcommand{\nphysa}{Nuclear Physics A}
\newcommand{\physrep}{Physics Reports}
\newcommand{\physscr}{Physica Scripta}
\newcommand{\planss}{Planetary Space Science}
\newcommand{\procspie}{Proceedings of the SPIE}

%% Abbreviations etc...

\newcommand{\astrolib}{\textsc{Astrolib}}
\newcommand{\ircamdr}{\textsc{Ircamdr}}
\newcommand{\ircam}{\textsc{Ircam}}
\newcommand{\daophot}{\textsc{Daophot}}
\newcommand{\ukirt}{\textsc{Ukirt}}
\newcommand{\hipp}{\textsc{Hipparcos}}
\newcommand{\baddata}{\textsc{Bad Data}}
\newcommand{\iraf}{\textsc{iraf}}
\newcommand{\isis}{\textsc{isis}}

\newcommand{\stare}{\textsc{Stare}}

%% Outer/inner scales...

\newcommand{\Lo}{\ensuremath{L_0}}
\newcommand{\lo}{\ensuremath{l_0}}

%% r_0, D/r_0...

\newcommand{\ro}{\ensuremath{r_0}}
\newcommand{\Dro}{\ensuremath{\frac{D}{\ro}}}
\newcommand{\roff}{\ensuremath{\ro(0.55\mu m)}}

\newcommand{\arcsec}{\mbox{\ensuremath{^{\prime\prime}}}}
\newcommand{\arcmin}{\mbox{\ensuremath{^{\prime}}}}
\newcommand{\arcpix}{\mbox{\ensuremath{^{\prime\prime}}/pixel}}
\newcommand{\mas}{\mbox{\ensuremath{mas}}}
\newcommand{\masyr}{\mbox{\ensuremath{mas\ yr^{-1}}}}

\newcommand{\cosa}{\mbox{\ensuremath{\cos\alpha}}}
\newcommand{\cosd}{\mbox{\ensuremath{\cos\delta}}}
\newcommand{\sina}{\mbox{\ensuremath{\sin\alpha}}}
\newcommand{\sind}{\mbox{\ensuremath{\sin\delta}}}
\newcommand{\seca}{\mbox{\ensuremath{\sec\alpha}}}
\newcommand{\secd}{\mbox{\ensuremath{\sec\delta}}}

\newcommand{\mR}{\ensuremath{m_R}}

\newcommand{\bolitem}[2]{
	\item \textbf{#1}

	#2
}
\newcommand{\nbolitem}[2]{
	\item #1

	#2
}

%% Maths stuff...

\renewcommand{\vec}[1]{\ensuremath{\underline{#1}}}	%% redefines
							%% an existing
							%% command...!
\newcommand{\curly}[1]{\ensuremath{\mathcal{#1}}}

%%NNets maths symbols...

\newcommand{\eqnspace}{\hspace{1.0cm}}
\newcommand{\oline}[1]{\ensuremath{\overline{#1}}}



%%my stuff from Mike's ddoc
\newcommand{\micron}{\mbox{\ensuremath{\,\mu}\rm{m}}}
\newcommand{\kel}{\ensuremath{\,\rm{K}}}
\newcommand{\power}[2]{\mbox{\ensuremath {#1 \times 10^{ #2 }}}}   % exp power
\newcommand{\Trot}{\ensuremath{T_{\rm{rot}}}}
\newcommand{\Lsolar}{\mbox{\ensuremath{\rm L_{\odot}}}}
\newcommand{\Lprime}{\ensuremath{L'}}
\newcommand{\Msolar}{\mbox{\ensuremath{\rm M_{\odot}}}}
\newcommand{\Zsolar}{\mbox{\ensuremath{\rm Z_{\odot}}}}
\newcommand{\kms}{\mbox{km\ensuremath{\rm{s}^{-1}}}}
\newcommand{\jy}{\ensuremath{\,\rm{Jy}}}
\newcommand{\ang}{\,\AA}
\newcommand{\ghz}{\,GHz}
\newcommand{\mhz}{\,MHz}

%%lick indices names
\newcommand{\mgb}{\ensuremath{\rm{Mg}_b}}
\newcommand{\mgone}{\ensuremath{\rm{Mg}_1}}
\newcommand{\mgtwo}{\ensuremath{\rm{Mg}_2}}
\newcommand{\ctwo}{\ensuremath{\rm{C}_2}}
\newcommand{\hdelta}{\ensuremath{{\rm H}\delta}}
\newcommand{\hdeltaa}{\ensuremath{{\rm H}\delta_A}}
\newcommand{\hdeltaf}{\ensuremath{{\rm H}\delta_F}}
\newcommand{\hgamma}{\ensuremath{{\rm H}\gamma}}
\newcommand{\hgammaa}{\ensuremath{{\rm H}\gamma_A}}
\newcommand{\hgammaf}{\ensuremath{{\rm H}\gamma_F}}


\newcommand{\degrees}{\mbox{\ensuremath{^{\circ}}}}
\newcommand{\electrons}{\ensuremath{\rm{e}^-}}

\newcommand{\hubsq}{\mbox{\ensuremath{h^{-2}}}}
\newcommand{\hub}{\mbox{\ensuremath{h^{-1}}}}

\newcommand{\logten}{\mbox{\ensuremath{{\rm log}_{10}}}}
\newcommand{\loge}{\mbox{\ensuremath{{\rm log}_e}}}

\newcommand{\intinf}{\ensuremath{\int_{- \infty}^{\infty}}}


%	MATRICES AND DETERMINANTS
%	=========================
%							UNIT ARRAY
\newcommand{\unitarray}{%
		\begin{array}{lll}
			 1	 &	 0	 &	 0	\\
			 0	 &	 1	 &	 0	\\
			 0	 &	 0	 &	 1	
		\end{array}%
		}
%							3 BY 3 ARRAY
\newcommand{\arrayyy}[1]{%
		\begin{array}{lll}
		 {#1}_{11}	 &	 {#1}_{12}	 &	 {#1}_{13}	\\
		 {#1}_{21}	 &	 {#1}_{22}	 &	 {#1}_{23}	\\
		 {#1}_{31}	 &	 {#1}_{32}	 &	 {#1}_{33}	
		\end{array}%
		}
%							2 BY 2 ARRAY
\newcommand{\arrayy}[1]{%
		\begin{array}{ll}
		 {#1}_{11}	 &	 {#1}_{12}	\\
		 {#1}_{21}	 &	 {#1}_{22}
		\end{array}%
		}
\newcommand{\unitm}{\left(\unitarray\right)}		% UNIT MATRIX
\newcommand{\unitd}{\left|\unitarray\right|}		% UNIT DET
\newcommand{\mattt}[1]{\left(\arrayyy{{#1}}\right)}	% 3*3  MATRIX
\newcommand{\matt}[1]{\left(\arrayy{{#1}}\right)}	% 2*2  MATRIX
\newcommand{\dettt}[1]{\left|\arrayyy{{#1}}\right|}	% 3*3  DET
\newcommand{\dett}[1]{\left|\arrayy{{#1}}\right|}	% 2*2  DET

\newcommand{\vrow}[1]{(#1_1,#1_2,#1_3)}			% ROW VECTOR
\newcommand{\myvrow}[2]{(#1_1,#1_2, . .\  #1_#2)}	% ROW VECTOR

\newcommand{\vcol}[1]{\left(\!%				% COL VECTOR
		\begin{array}{l}
			{#1}_{1}	\\
			{#1}_{2}	\\
			{#1}_{3}
		\end{array}%
		\!\right)}
\newcommand{\myvcol}[2]{\left(\!%				% COL VECTOR
		\begin{array}{l}
			{#1}_{1}	\\
			{#1}_{2}	\\
			. \\
			. \\
			{#1}_{#2}
		\end{array}%
		\!\right)}

\newcommand{\mydesign}[3]{\left(\!%			% Design matrix
		\begin{array}{lllll}
			{#1}_{11} & {#1}_{12} & . & . & {#1}_{1#3} \\
			{#1}_{21} & {#1}_{22} & . & . & {#1}_{2#3} \\
			. & & & & . \\
			. & & & & . \\
			{#1}_{#21} & {#1}_{#22} & . & . & {#1}_{#2#3} \\
		\end{array}%
		\!\right)}

\def\dmatrix#1{\left|\matrix{#1}\right|}

%		 CALCULUS
%		========
\newcommand{\de}{\,{\rm d}}				% d for dx etc.
\newcommand{\dd}[2]{\frac{\de#1}{\de#2}}		% derivative df/dx
\newcommand{\pp}[2]{\frac{\partial#1}{\partial#2}}	% partial df/dx
\newcommand{\ff}[2]{\frac{\delta#1}{\delta#2}}		% functional df/dx
\newcommand{\ddt}[2]{\frac{\de^{2}\! #1}{\de#2^{2}}}	% 2nd der: d2f/dx2
\newcommand{\ppt}[2]					% 2nd partial der.
	{\frac{\partial^{2}\! #1}{\partial#2^{2}}}
\newcommand{\ppm}[3]						% 2nd partial mixed
	{\frac{\partial^{2}\! #1}{\partial#2\partial#3}}	
\newcommand{\fft}[2]
	{\frac{\delta^{2}\! #1}{\delta#2^{2}}}		% 2nd functional der.
\newcommand{\ddn}[3]
	{\frac{\de^{#1}\! #2}{\de#3^{#1}}}		% nth derivative
\newcommand{\ppn}[3]
	{\frac{\partial^{#1}\! #2}{\partial#3^{#1}}}	% nth partial der.
\newcommand{\ffn}[3]
	{\frac{\delta^{#1}\!#2}{\delta#3^{#1}}}		% nth functional der.

%% Reference yet to find...

%%\newcommand{\findcite}[1]{[!!! Find ref to: #1!!!]}
%%\newcommand{\findref}[1]{[!!! Need to cross--reference to: #1!!!]}

\newcommand{\todo}[1]{
\vspace{0.5in}
\textbf{Todo: #1}
\vspace{0.5in}
}

\newcommand{\next}[1]{
\vspace{2in}
\begin{large}
\textbf{Next: #1}
\end{large}
\vspace{2in}
}

\newcommand{\grumpy}				% draws a grumpy face
{\mbox{\begin{picture}(20,20)(0,7)
\put(10,10){\circle{15}}
\put(12.5,12.5){\circle*{2}}
\put(7.5,12.5){\circle*{2}}
\put(6,1){\shortstack[l]{$\mbox{}^{\frown}$}}
\end{picture}}}

\newcommand{\happy}				% draws a happy face
{\mbox{\begin{picture}(20,20)(0,7)
\put(10,10){\circle{15}}
\put(12.5,12.5){\circle*{2}}
\put(7.5,12.5){\circle*{2}}
\put(6,1){\shortstack[l]{$\mbox{}^{\smile}$}}
\end{picture}}}

%% The keys to the graphs...

\newcommand{\xxx}{		%% Cross...
	$\!\!\!\!\!$
	\setlength{\unitlength}{1pt}
	\mbox{
	\begin{picture}(16,4)(0,0)
	\put(8,2){\makebox(0,0){$\times$}}
	\put(0,2){\line(1,0){16}}
	\end{picture}
	\setlength{\unitlength}{1pt}
	$\!\!\!\!\!$
	}
}

\newcommand{\xxxns}{		%% As above, no symbol...
	$\!\!\!\!\!$
	\setlength{\unitlength}{1pt}
	\mbox{
	\begin{picture}(16,4)(0,0)
	\put(0,2){\line(1,0){16}}
	\end{picture}
	\setlength{\unitlength}{1pt}
	$\!\!\!\!\!$
	}
}

\newcommand{\ddd}{		%% Diamond...
	$\!\!\!\!\!$
	\setlength{\unitlength}{1pt}
	\mbox{
	\begin{picture}(16,4)(0,0)
	\put(8,2){\makebox(0,0){$\diamond$}}
	\put(0,2){\line(1,0){4}}
	\put(6,2){\line(1,0){4}}
	\put(12,2){\line(1,0){4}}
	\end{picture}
	\setlength{\unitlength}{1pt}
	$\!\!\!\!\!$
	}
}

\newcommand{\dddns}{		%% As above, no symbol...
	$\!\!\!\!\!$
	\setlength{\unitlength}{1pt}
	\mbox{
	\begin{picture}(16,4)(0,0)
	\put(0,2){\line(1,0){4}}
	\put(6,2){\line(1,0){4}}
	\put(12,2){\line(1,0){4}}
	\end{picture}
	\setlength{\unitlength}{1pt}
	$\!\!\!\!\!$
	}
}

\newcommand{\ttt}{		%% Triangle...
	$\!\!\!\!\!$
	\setlength{\unitlength}{1pt}
	\mbox{
	\begin{picture}(16,4)(0,0)
	\put(8,2){\makebox(0,0){\begin{tiny}$\triangle$\end{tiny}}}
	\put(0,2){\line(1,0){3}}
	\put(4,2){\line(1,0){1}}
	\put(6,2){\line(1,0){3}}
	\put(10,2){\line(1,0){1}}
	\put(12,2){\line(1,0){3}}
	\end{picture}
	\setlength{\unitlength}{1pt}
	$\!\!\!\!\!$
	}
}

\newcommand{\key}[3]{\\\mbox{\xxx} #1,\\\mbox{\ddd} #2 and\\\mbox{\ttt} #3}
\newcommand{\keydef}{\key{no prediction}{neural network}{linear predictor}}
\newcommand{\keydefs}{\key{no prediction}{neural networks}{linear predictors}}

%% A figure...

\newcommand{\fig}[4]{
	\begin{figure}
	\setlength{\epsfxsize}{0.8 \textwidth}
	\centerline{\epsfbox{#1}}
	\caption[#2]{ \label{#4} #3}
	\vspace{18pt}
	\end{figure}
}

%% A scaleable figure...

\newcommand{\scfig}[5]{
	\begin{figure}
	\setlength{\epsfxsize}{#1 \textwidth}
	\centerline{\epsfbox{#2}}
	\caption[#3]{ \label{#5} #4}
	\vspace{18pt}
	\end{figure}
}

%% Fitting 2 graphs into one figure on one page...

\newcommand{\pfig}[7]{
	\begin{figure}
	\begin{center}
	\begin{tabular}{c}
	\setlength{\epsfxsize}{0.8 \textwidth} \epsfbox{#1} \\
	(a) \begin{minipage}[t]{0.75 \textwidth} #2 \end{minipage} \\
	\\
	\setlength{\epsfxsize}{0.8 \textwidth} \epsfbox{#3} \\
	(b) \begin{minipage}[t]{0.75 \textwidth} #4 \end{minipage} \\
	\end{tabular}
	\end{center}
	\caption[#5]{\label{#7} #6}
	\end{figure}
}

\newcommand{\scpfig}[9]{
	\begin{figure}
	\begin{center}
	\begin{tabular}{c}
	\setlength{\epsfxsize}{#8 \textwidth} \epsfbox{#1} \\
	(a) \begin{minipage}[t]{0.75 \textwidth} #2 \end{minipage} \\
	\\
	\setlength{\epsfxsize}{#9 \textwidth} \epsfbox{#3} \\
	(b) \begin{minipage}[t]{0.75 \textwidth} #4 \end{minipage} \\
	\end{tabular}
	\end{center}
	\caption[#5]{\label{#7} #6}
	\end{figure}
}

%% Fitting another two graphs onto a second page, retaining figure
%% number of the previous...

\newcommand{\ppfig}[7]{
	\begin{figure}
	\addtocounter{figure}{-1}
	\begin{center}
	\begin{tabular}{c}
	\setlength{\epsfxsize}{0.8 \textwidth} \epsfbox{#1} \\
	(c) \begin{minipage}[t]{0.75 \textwidth} #2 \end{minipage} \\
	\\
	\setlength{\epsfxsize}{0.8 \textwidth} \epsfbox{#3} \\
	(d) \begin{minipage}[t]{0.75 \textwidth} #4 \end{minipage} \\
	\end{tabular}
	\end{center}
	\caption[$\!\!$(\emph{cont.}) #5]
		{\label{#7_cont} $\!$(\emph{cont.}): #6}
	\end{figure}
}

%%
%% End of file...
%%